\section{Stability Analysis}

\begin{purposebox}
\textbf{Purpose:} Demonstrate understanding of Lyapunov stability, existence/uniqueness of steady states, finite-time convergence, and sliding mode analysis---the theoretical backbone proving the defense mechanism works.

\textbf{When complete:} You can independently walk through the paper's stability proofs and explain why the defended tracking control NN is guaranteed to work correctly.
\end{purposebox}

%──────────────────────────────────────────────────
\subsection{Lyapunov Stability Theory}
%──────────────────────────────────────────────────

\subsubsection*{Concept Check}

\question{Lyapunov's Second Method}
State the two conditions required for Lyapunov stability:
\begin{enumerate}[nosep]
  \item Condition on $V(x)$: ?
  \item Condition on $\Delta V(k) = V(k+1) - V(k)$: ?
\end{enumerate}
In your own words, explain the intuition: why do these two conditions guarantee stability?

\begin{answerbox}\end{answerbox}

\question{Lyapunov Function Construction}
The paper uses $W(k) = |x(k)|$ as the Lyapunov candidate function. Verify that it satisfies the first condition (positive definite): Is $W(x) > 0$ for $x \neq 0$ and $W(0) = 0$?

\begin{answerbox}\end{answerbox}

\question{Global Asymptotic Stability}
Define global asymptotic stability. What is the difference between:
\begin{itemize}[nosep]
  \item Stability (trajectories stay close)
  \item Asymptotic stability (trajectories converge)
  \item Global asymptotic stability (converge from \emph{any} initial condition)
\end{itemize}
Why is \emph{global} stability important for this project's defense guarantees?

\begin{answerbox}\end{answerbox}

\question{Radial Unboundedness}
What does $V(x) \to \infty$ as $\|x\| \to \infty$ mean? Why is this condition needed for \emph{global} (not just local) stability?

\begin{answerbox}\end{answerbox}

%──────────────────────────────────────────────────
\subsection{Existence \& Uniqueness of Steady States}
%──────────────────────────────────────────────────

\subsubsection*{Concept Check}

\question{Equilibrium Analysis}
At a steady state, $x(k+1) = x(k)$, so the change is zero. Write the equilibrium condition from the tracking equation:
\[
  \frac{x(k+1) - x(k)}{\tau} = h(x)\,S(x) = 0
\]
What are the two ways this product can equal zero? Which corresponds to the correct steady state?

\begin{answerbox}\end{answerbox}

\question{Why Uniqueness Matters}
Suppose the system had multiple steady states. Explain in your own words why this would be problematic for:
\begin{enumerate}[nosep, label=(\alph*)]
  \item Normal operation (tracking a reference)
  \item Operation under adversarial attack
\end{enumerate}

\begin{answerbox}\end{answerbox}

\question{Connecting to Defense}
If an adversary perturbs the system, the steady state analysis guarantees the system will still converge to the \emph{correct} steady state (not a false one). Explain why having a unique steady state is the foundation of the defense mechanism.

\begin{answerbox}\end{answerbox}

%──────────────────────────────────────────────────
\subsection{Finite-Time Convergence}
%──────────────────────────────────────────────────

\subsubsection*{Concept Check}

\question{Lemma~2 Statement}
State Lemma~2 in your own words: trajectories converge to a $\delta$-neighborhood of the steady state in at most $m$ steps, where $m \leq (x(1) - \delta)/\sigma + 1$ and $\sigma = \tau\,h(\delta)$.

What do each of these quantities represent?
\begin{itemize}[nosep]
  \item $\delta$: ?
  \item $\sigma$: ?
  \item $m$: ?
\end{itemize}

\begin{answerbox}\end{answerbox}

\question{Convergence Bound Derivation}
The proof uses $W(x(m)) \leq W(x(1)) - \sigma(m-1)$. Explain the logic:
\begin{itemize}[nosep]
  \item Why does $W$ decrease by at least $\sigma$ each step?
  \item How does setting $W(x(m)) \geq 0$ give you the upper bound on $m$?
\end{itemize}

\begin{largeanswer}\end{largeanswer}

\question{Invariant Set Property}
Once the trajectory enters the $\delta$-neighborhood, it stays there (the set $\Phi$ is invariant). Why is this property important? What would happen without it?

\begin{answerbox}\end{answerbox}

%──────────────────────────────────────────────────
\subsection{Sliding Mode Analysis}
%──────────────────────────────────────────────────

\subsubsection*{Concept Check}

\question{Sliding Mode Conditions}
State the two sliding mode conditions:
\begin{itemize}[nosep]
  \item $\Delta x(k) > 0$ when $x(k) \to -0$
  \item $\Delta x(k) < 0$ when $x(k) \to +0$
\end{itemize}
Draw a simple diagram or describe in words what this means: the trajectory oscillates around $x = 0$.

\begin{largeanswer}\end{largeanswer}

\question{Chattering Phenomenon}
What is chattering? Why does it occur in discrete-time systems with signum switching? Is chattering a problem or a feature in this context?

\begin{answerbox}\end{answerbox}

\question{Full Stability Story}
Put it all together. Trace the complete stability narrative:
\begin{enumerate}[nosep]
  \item System starts at arbitrary initial condition $x(0)$
  \item[\textrm{2--5.}] \emph{What happens next? Fill in each phase until steady state.}
\end{enumerate}

\begin{largeanswer}\end{largeanswer}

%──────────────────────────────────────────────────
\subsection*{Self-Assessment Checklist}
%──────────────────────────────────────────────────

\begin{itemize}[leftmargin=2em, label=$\square$]
  \item I can state Lyapunov's second method and verify a candidate function
  \item I can explain global asymptotic stability and radial unboundedness
  \item I can derive the equilibrium condition and explain uniqueness
  \item I can state the finite-time convergence bound and explain the proof logic
  \item I can describe sliding mode behavior and chattering
  \item I can tell the complete stability story from initial condition to steady state
\end{itemize}

\begin{notesbox}\end{notesbox}
